% GENERATED FILE. Edit canonical lessons, then run npm run generate:books.
% canonical-source-hash: fnv1a64:d788a402b1c0a64b
% canonical-lessons: BN-C361-shikkha, BN-C361-bidya, BN-C361-odhyapok, BN-C361-digri, BN-C361-porano-teach, BN-R361-review-354-361

\chapter{Learning and Teaching}
\label{ch:bn-learning-and-teaching}
\hlchaptermodality{361}

\begin{chapteropening}
\textbf{By the end of this chapter:} \emph{I can talk about education and who teaches what.}

Complete the last lesson of chapter 361: I can talk about education and who teaches what.
\end{chapteropening}

\section[\texorpdfstring{\bn{শিক্ষা} (shikkhā)}{শিক্ষা (shikkhā)}]{\bn{শিক্ষা} (shikkhā) — education}

\label{lesson:BN-C361-shikkha}



\begin{quote}
Before the new one: say the Bengali for in attendance, then the Bengali for to address a crowd.
\end{quote}

\subsection*{You'll want to know: \bn{শিক্ষা}}
\textbf{\bn{শিক্ষা}} — \emph{shikkhā} — \textquotedblleft{}education\textquotedblright{}.

\begin{quote}
\bn{শিক্ষা} \bn{সবার} \bn{অধিকার।}
\end{quote}

\begin{quote}
\emph{shikkhā sôbār ôdhikār.}
\end{quote}

\begin{quote}
Education is everyone's right.
\end{quote}

\begin{cousinweb}
\textbf{\bn{শিক্ষা}} (\emph{shikkhā}) is Sanskrit \emph{śikṣā}, learning or study (Monier-Williams). The everyday verb \bn{শেখা}, to learn, comes from the same root by the spoken road (CDIAL 12430 \emph{śíkṣatē} lists Bengali sikhā).
\end{cousinweb}

\begin{grammarlens}[title={What you've built}]
One more word for giving reasons and opinions.
\end{grammarlens}

\subsection*{Your turn}
\begin{itemize}
\raggedright
  \item \emph{Say it:} \emph{shikkhā}
  \item \emph{Say it:} \emph{shikkhā}, once more
  \item \emph{Say it:} \emph{shikkhā sôbār ôdhikār.} — Education is everyone's right.
  \item \emph{Recall:} say the Bengali for in attendance, then the Bengali for to address a crowd, then say \emph{shikkhā} again
\end{itemize}

\begin{tcolorbox}[breakable,colback=teal!4,colframe=teal!35!black,title={Before you move on}]
What does \bn{শিক্ষা} mean? (\textquotedblleft{}Education\textquotedblright{}.) Say it once more.
\end{tcolorbox}
\section[\texorpdfstring{\bn{বিদ্যা} (bidyā)}{বিদ্যা (bidyā)}]{\bn{বিদ্যা} (bidyā) — learning; knowledge}

\label{lesson:BN-C361-bidya}



\begin{quote}
Before the new one: say the Bengali for to address a crowd, then the Bengali for education.
\end{quote}

\subsection*{You'll want to know: \bn{বিদ্যা}}
\textbf{\bn{বিদ্যা}} — \emph{bidyā} — \textquotedblleft{}learning; knowledge\textquotedblright{}.

\begin{quote}
\bn{বিদ্যা} \bn{কেউ} \bn{চুরি} \bn{করতে} \bn{পারে} \bn{না।}
\end{quote}

\begin{quote}
\emph{bidyā keu churi kôrte pāre nā.}
\end{quote}

\begin{quote}
No one can steal your learning.
\end{quote}

\begin{cousinweb}
\textbf{\bn{বিদ্যা}} (\emph{bidyā}) is Sanskrit \emph{vidyā}, knowledge or learning (Monier-Williams). It is the learning a person carries, as in \bn{বিদ্যালয়}, a school, 'the house of learning'.
\end{cousinweb}

\begin{grammarlens}[title={What you've built}]
One more word for giving reasons and opinions.
\end{grammarlens}

\subsection*{Your turn}
\begin{itemize}
\raggedright
  \item \emph{Say it:} \emph{bidyā}
  \item \emph{Say it:} \emph{bidyā}, once more
  \item \emph{Say it:} \emph{bidyā keu churi kôrte pāre nā.} — No one can steal your learning.
  \item \emph{Recall:} say the Bengali for to address a crowd, then the Bengali for education, then say \emph{bidyā} again
\end{itemize}

\begin{tcolorbox}[breakable,colback=teal!4,colframe=teal!35!black,title={Before you move on}]
What does \bn{বিদ্যা} mean? (\textquotedblleft{}Learning; knowledge\textquotedblright{}.) Say it once more.
\end{tcolorbox}
\section[\texorpdfstring{\bn{অধ্যাপক} (ôdhyāpôk)}{অধ্যাপক (ôdhyāpôk)}]{\bn{অধ্যাপক} (ôdhyāpôk) — a professor; a college lecturer}

\label{lesson:BN-C361-odhyapok}



\begin{quote}
Before the new one: say the Bengali for education, then the Bengali for learning.
\end{quote}

\subsection*{You'll want to know: \bn{অধ্যাপক}}
\textbf{\bn{অধ্যাপক}} — \emph{ôdhyāpôk} — \textquotedblleft{}a professor; a college lecturer\textquotedblright{}.

\begin{quote}
\bn{মীরার} \bn{বাবা} \bn{কলেজে} \bn{অধ্যাপক।}
\end{quote}

\begin{quote}
\emph{mīrār bābā kôleje ôdhyāpôk.}
\end{quote}

\begin{quote}
Mira's father is a lecturer at the college.
\end{quote}

\begin{cousinweb}
\textbf{\bn{অধ্যাপক}} (\emph{ôdhyāpôk}) is Sanskrit \emph{adhyāpaka}, a teacher (Monier-Williams). Bengali keeps it for those who teach at a college or a university.
\end{cousinweb}

\begin{grammarlens}[title={What you've built}]
One more word for giving reasons and opinions.
\end{grammarlens}

\subsection*{Your turn}
\begin{itemize}
\raggedright
  \item \emph{Say it:} \emph{ôdhyāpôk}
  \item \emph{Say it:} \emph{ôdhyāpôk}, once more
  \item \emph{Say it:} \emph{mīrār bābā kôleje ôdhyāpôk.} — Mira's father is a lecturer at the college.
  \item \emph{Recall:} say the Bengali for education, then the Bengali for learning, then say \emph{ôdhyāpôk} again
\end{itemize}

\begin{tcolorbox}[breakable,colback=teal!4,colframe=teal!35!black,title={Before you move on}]
What does \bn{অধ্যাপক} mean? (\textquotedblleft{}A professor; a college lecturer\textquotedblright{}.) Say it once more.
\end{tcolorbox}
\section[\texorpdfstring{\bn{ডিগ্রি} (ḍigri)}{ডিগ্রি (ḍigri)}]{\bn{ডিগ্রি} (ḍigri) — a university degree}

\label{lesson:BN-C361-digri}



\begin{quote}
Before the new one: say the Bengali for learning, then the Bengali for a professor.
\end{quote}

\subsection*{You'll want to know: \bn{ডিগ্রি}}
\textbf{\bn{ডিগ্রি}} — \emph{ḍigri} — \textquotedblleft{}a university degree\textquotedblright{}.

\begin{quote}
\bn{অরুণ} \bn{বিজ্ঞানে} \bn{ডিগ্রি} \bn{পেল।}
\end{quote}

\begin{quote}
\emph{ôrun bigyāne ḍigri pelô.}
\end{quote}

\begin{quote}
Arun got a degree in science.
\end{quote}

\begin{cousinweb}
\textbf{\bn{ডিগ্রি}} (\emph{ḍigri}) is borrowed from English \emph{degree}, with the English d heard as the Bengali ḍ.
\end{cousinweb}

\begin{grammarlens}[title={What you've built}]
One more word for giving reasons and opinions.
\end{grammarlens}

\subsection*{Your turn}
\begin{itemize}
\raggedright
  \item \emph{Say it:} \emph{ḍigri}
  \item \emph{Say it:} \emph{ḍigri}, once more
  \item \emph{Say it:} \emph{ôrun bigyāne ḍigri pelô.} — Arun got a degree in science.
  \item \emph{Recall:} say the Bengali for learning, then the Bengali for a professor, then say \emph{ḍigri} again
\end{itemize}

\begin{tcolorbox}[breakable,colback=teal!4,colframe=teal!35!black,title={Before you move on}]
What does \bn{ডিগ্রি} mean? (\textquotedblleft{}A university degree\textquotedblright{}.) Say it once more.
\end{tcolorbox}
\section[\texorpdfstring{\bn{পড়ানো} (pôṛāno)}{পড়ানো (pôṛāno)}]{\bn{পড়ানো} (pôṛāno) — to teach (a subject); to tutor}

\label{lesson:BN-C361-porano-teach}



\begin{quote}
Before the new one: say the Bengali for a professor, then the Bengali for a university degree.
\end{quote}

\subsection*{You'll want to know: \bn{পড়ানো}}
\textbf{\bn{পড়ানো}} — \emph{pôṛāno} — \textquotedblleft{}to teach (a subject); to tutor\textquotedblright{}.

\begin{quote}
\bn{মীরা} \bn{সন্ধ্যায়} \bn{শিশুদের} \bn{ইংরেজি} \bn{পড়ায়।}
\end{quote}

\begin{quote}
\emph{mīrā sôndhyāy shishuder ingreji pôṛāy.}
\end{quote}

\begin{quote}
Mira teaches children English in the evenings.
\end{quote}

Past (he/she …): \bn{তিনি} \bn{অনেক} \bn{বছর} \bn{ইতিহাস} \bn{পড়ালেন}, \emph{tini ônek bôchhôr itihās pôṛālen}.

\begin{cousinweb}
\textbf{\bn{পড়ানো}} (\emph{pôṛāno}) is the causative of \bn{পড়া} 'to read, to study' (CDIAL 7712 \emph{páṭhati} lists Bengali paṛā): to make someone study. \bn{শেখানো} teaches a skill; \bn{পড়ানো} teaches a subject from books.
\end{cousinweb}

\begin{grammarlens}[title={What you've built}]
One more word for giving reasons and opinions.
\end{grammarlens}

\subsection*{Your turn}
\begin{itemize}
\raggedright
  \item \emph{Say it:} \emph{pôṛāno}
  \item \emph{Say it:} \emph{pôṛāno}, once more
  \item \emph{Say it:} \emph{mīrā sôndhyāy shishuder ingreji pôṛāy.} — Mira teaches children English in the evenings.
  \item \emph{Recall:} say the Bengali for a professor, then the Bengali for a university degree, then say \emph{pôṛāno} again
\end{itemize}

\begin{tcolorbox}[breakable,colback=teal!4,colframe=teal!35!black,title={Before you move on}]
What does \bn{পড়ানো} mean? (\textquotedblleft{}To teach (a subject); to tutor\textquotedblright{}.) Say it once more.
\end{tcolorbox}
\section[(dialogue)]{Looking back: from Finding Out to Learning and Teaching}

\label{lesson:BN-R361-review-354-361}



\begin{quote}
Say the Bengali for research and an experiment.
\end{quote}

\subsection*{Your turn}
Say these aloud:

\begin{itemize}
\raggedright
  \item research, an experiment, a theory, to discover, to look into
  \item statistics, a sample, a survey, to collect, to analyse
  \item a percentage, a ratio, to calculate, a yardstick, to assess
  \item an essay, a subject, a remark, a conclusion of a text, a notion
  \item a likeness, a difference, a link, ordinary, to mix up
  \item an excuse, a pretence, a demanded explanation, to sidestep, acceptable
  \item a lecture, a speaker, a chairperson, in attendance, to address a crowd
  \item education, learning, a professor, a university degree, to tutor
\end{itemize}

\begin{tcolorbox}[breakable,colback=teal!4,colframe=teal!35!black,title={Before you move on}]
Research? (\textbf{\bn{গবেষণা}}, \emph{gôbeshônā}.) An experiment? (\textbf{\bn{পরীক্ষা}}, \emph{pôrīkkhā}.) A theory? (\textbf{\bn{তত্ত্ব}}, \emph{tôttvô}.) To discover? (\textbf{\bn{আবিষ্কার} \bn{করা}}, \emph{ābishkār kôrā}.) To look into? (\textbf{\bn{অনুসন্ধান} \bn{করা}}, \emph{ônusôndhān kôrā}.) Statistics? (\textbf{\bn{পরিসংখ্যান}}, \emph{pôrisôngkhyān}.) A sample? (\textbf{\bn{নমুনা}}, \emph{nômunā}.) A survey? (\textbf{\bn{জরিপ}}, \emph{jôrip}.) To collect? (\textbf{\bn{সংগ্রহ} \bn{করা}}, \emph{sônggrôhô kôrā}.) To analyse? (\textbf{\bn{বিশ্লেষণ} \bn{করা}}, \emph{bishleshôn kôrā}.) A percentage? (\textbf{\bn{শতাংশ}}, \emph{shôtāngshô}.) A ratio? (\textbf{\bn{অনুপাত}}, \emph{ônupāt}.) To calculate? (\textbf{\bn{হিসেব} \bn{করা}}, \emph{hiseb kôrā}.) A yardstick? (\textbf{\bn{মাপকাঠি}}, \emph{māpkāṭhi}.) To assess? (\textbf{\bn{মূল্যায়ন} \bn{করা}}, \emph{mūlyāyôn kôrā}.) An essay? (\textbf{\bn{প্রবন্ধ}}, \emph{prôbôndhô}.) A subject? (\textbf{\bn{বিষয়}}, \emph{bishôy}.) A remark? (\textbf{\bn{মন্তব্য}}, \emph{môntôbyô}.) A conclusion of a text? (\textbf{\bn{উপসংহার}}, \emph{upôsônghār}.) A notion? (\textbf{\bn{ধারণা}}, \emph{dhārônā}.) A likeness? (\textbf{\bn{মিল}}, \emph{mil}.) A difference? (\textbf{\bn{পার্থক্য}}, \emph{pārthôkyô}.) A link? (\textbf{\bn{যোগসূত্র}}, \emph{jogsūtrô}.) Ordinary? (\textbf{\bn{সাধারণ}}, \emph{sādhārôn}.) To mix up? (\textbf{\bn{গুলিয়ে} \bn{ফেলা}}, \emph{guliye phelā}.) An excuse? (\textbf{\bn{অজুহাত}}, \emph{ôjuhāt}.) A pretence? (\textbf{\bn{ভান}}, \emph{bhān}.) A demanded explanation? (\textbf{\bn{কৈফিয়ত}}, \emph{koiphiyôt}.) To sidestep? (\textbf{\bn{পাশ} \bn{কাটানো}}, \emph{pāsh kāṭāno}.) Acceptable? (\textbf{\bn{গ্রহণযোগ্য}}, \emph{grôhônjogyô}.) A lecture? (\textbf{\bn{বক্তৃতা}}, \emph{bôktritā}.) A speaker? (\textbf{\bn{বক্তা}}, \emph{bôktā}.) A chairperson? (\textbf{\bn{সভাপতি}}, \emph{sôbhāpôti}.) In attendance? (\textbf{\bn{উপস্থিত}}, \emph{upôsthitô}.) To address a crowd? (\textbf{\bn{সম্বোধন} \bn{করা}}, \emph{sômbodhôn kôrā}.) Education? (\textbf{\bn{শিক্ষা}}, \emph{shikkhā}.) Learning? (\textbf{\bn{বিদ্যা}}, \emph{bidyā}.) A professor? (\textbf{\bn{অধ্যাপক}}, \emph{ôdhyāpôk}.) A university degree? (\textbf{\bn{ডিগ্রি}}, \emph{ḍigri}.) To tutor? (\textbf{\bn{পড়ানো}}, \emph{pôṛāno}.)
\end{tcolorbox}
